\section{Introduction}
\label{sec:intro}

Routine UAV inspection is a robotics perception problem with an unusually large gap between what the camera sees and what an operator must decide. A wide aerial frame may contain roads, houses, waterways, vehicles, and work sites, but the actionable output is an \emph{operational event}: a fire lane is obstructed, waste floats on water rather than lying on the bank, a rural building is under construction, or stored material is actually debris. These labels compose an entity, state, relation, and scene context. The evidence is often a small region, while the correct interpretation depends on the full scene.

Recent low-altitude vision-language resources have raised the scale bar. UAVBench/UAVIT-1M evaluates ten image/region tasks and supplies more than a million instructions~\cite{zhan2026uavbench}; SpatialSky-Bench covers 13 UAV spatial-intelligence subcategories with a million-sample training resource~\cite{zhang2026spatialsky}; UAVReason unifies spatial/temporal reasoning with aerial generation in simulation~\cite{sun2026uavreason}; and AirNav targets real-world UAV vision-language navigation~\cite{cai2026airnav}. Consequently, ``UAV imagery plus VLM fine-tuning'' is no longer a sufficient contribution. What remains underrepresented is operationally grounded, fine-grained event understanding in repeated real inspection flights, where definitions are policy-facing, classes form explicit confusion sets, and the system must know when evidence is insufficient.

Our data expose this problem in a form that generic benchmarks do not. Each visible annotation links a wide \emph{context view} to a local \emph{evidence view}; repeated exact contexts may contain several crops and several valid events. The catalog also contains closely related concepts such as floating garbage, waterside garbage, unclean water, and algal bloom; or rural construction, accumulated building materials, and construction waste. These cases break object-name shortcuts and make naive pair-level splitting invalid: crops from the same scene, near-duplicate frames, and repeated flights can leak nearly identical evidence across train and test.

\begin{figure*}[t]
\centering
\begin{minipage}[t]{0.32\textwidth}
  \centering
  \includegraphics[width=0.61\linewidth,height=2.35cm,keepaspectratio]{street_vendor_context.jpg}\hfill
  \includegraphics[width=0.37\linewidth,height=2.35cm,keepaspectratio]{street_vendor_evidence.jpg}
  \par\vspace{2pt}\footnotesize\textbf{(a) Street vendor.} Local identity plus roadway occupation.
\end{minipage}\hfill
\begin{minipage}[t]{0.32\textwidth}
  \centering
  \includegraphics[width=0.61\linewidth,height=2.35cm,keepaspectratio]{floating_garbage_context.jpg}\hfill
  \includegraphics[width=0.37\linewidth,height=2.35cm,keepaspectratio]{floating_garbage_evidence.jpg}
  \par\vspace{2pt}\footnotesize\textbf{(b) Floating garbage.} Debris must lie on water, not the bank.
\end{minipage}\hfill
\begin{minipage}[t]{0.32\textwidth}
  \centering
  \includegraphics[width=0.61\linewidth,height=2.35cm,keepaspectratio]{construction_waste_context.jpg}\hfill
  \includegraphics[width=0.37\linewidth,height=2.35cm,keepaspectratio]{construction_waste_evidence.jpg}
  \par\vspace{2pt}\footnotesize\textbf{(c) Construction waste.} Material, site context, and state interact.
\end{minipage}
\caption{Each annotation links a global context (left) to local evidence (right); exact-context grouping further joins multiple crops and events from the same scene. Operational labels require factor composition, not object naming. Images are internal examples pending release clearance.}
\label{fig:teaser}
\end{figure*}


We reorganize the project around three ideas. First, a five-factor \eventgraph represents every label as a composition of domain, entity, state, relation, and context, with explicit neighbor edges for confusable concepts. Second, a \samplegraph preserves the empirical many-to-many relation among a wide context, its evidence crops, and its event set; this supports both pair recognition and context-level set prediction. Third, the benchmark is support-gated and leakage-aware, while \method contrasts each positive event against graph neighbors and aggregates a variable-size evidence set without imposing crop order.

This paper makes three contributions:
\begin{itemize}
    \item We introduce \dataset, a real low-altitude inspection benchmark with 9,093 reconciled pair annotations over 7,279 exact-hash-unique contexts and a 67-concept operational catalog. We publish audit gates, content-group identities, and versioned adjudication ledgers instead of treating pair rows or raw records as independent images.
    \item We contribute an \eventgraph, a many-to-many \samplegraph, evidence cards, and leakage-resistant tasks spanning pair recognition, multi-evidence event-set prediction, hard negatives, open-vocabulary transfer, unseen locations, selective prediction, and explicit tests of whether a second view is necessary.
    \item We propose \method, a Qwen3-VL-based parameter-efficient baseline that combines context and evidence with taxonomy-neighbor ranking and counterfactual rejection; its deterministic set extension exposes which crop supports each predicted event.
\end{itemize}

The manuscript is intentionally result-honest: verified snapshot statistics are included, whereas unresolved split counts and measurements remain \tbd. The submission version must replace every placeholder with outputs from the frozen protocol in \cref{sec:experiments}.
